% !TEX program = xelatex
\documentclass[12pt,a4paper]{article}

%%
% CUMCM 数学建模竞赛论文写作培训模板（2026规范+AI规定内嵌版）
% 用法：本文件既可作为培训讲义，也可作为参赛论文写作母版。
% 正式参赛时：删除所有“【培训提示】”文字、示例性内容，并替换为本队实际内容。
%%

\usepackage{geometry}
\geometry{a4paper,left=2.5cm,right=2.5cm,top=2.5cm,bottom=2.5cm}
\usepackage{ctex}
\usepackage{sectsty}
\sectionfont{\centering}
\renewcommand\thesection{\chinese{section}、}
\renewcommand\thesubsection{\arabic{section}.\arabic{subsection}}
\usepackage[dvipsnames]{xcolor}
\usepackage{booktabs}
\usepackage{makecell}
\usepackage{amsfonts,amsmath,amssymb,mathtools}
\usepackage{multirow}
\usepackage{graphicx}
\usepackage{subcaption}
\usepackage{float}
\usepackage{enumitem}
\usepackage{array}
\usepackage{longtable}
\usepackage{listings}
\usepackage{fancyhdr}
\usepackage{hyperref}
\usepackage{url}

\hypersetup{
  colorlinks=true,
  linkcolor=black,
  urlcolor=blue,
  citecolor=black,
  pdfstartview=Fit,
  breaklinks=true
}

\lstset{
  frame=tb,
  aboveskip=3mm,
  belowskip=3mm,
  showstringspaces=false,
  columns=flexible,
  framerule=1pt,
  rulecolor=\color{gray!35},
  backgroundcolor=\color{gray!5},
  basicstyle={\small\ttfamily},
  numbers=left,
  numberstyle=\tiny\color{gray},
  keywordstyle=\color{blue},
  commentstyle=\color{ForestGreen},
  breaklines=true,
  breakatwhitespace=true,
  tabsize=3
}
\renewcommand{\lstlistingname}{源程序}

\newcommand{\tip}[1]{\par\noindent\textbf{【提示】} #1\par}
\newcommand{\warning}[1]{\par\noindent\textbf{【高频失分点】} #1\par}
\newcommand{\good}[1]{\par\noindent\textbf{【推荐写法】} #1\par}
\newcommand{\bad}[1]{\par\noindent\textbf{【不推荐写法】} #1\par}
\newcommand{\core}[1]{\boxed{\text{#1}}}

\fancyhf{}
\cfoot{\thepage}
\renewcommand{\headrulewidth}{0pt}
\pagestyle{fancy}


% ============================================================
% 【新增·写作辅助宏】只用于“提示”，不产生论文正文内容
% ============================================================
\usepackage[font=small,labelsep=quad]{caption}

\newcommand{\prompt}[1]{\par\noindent\textbf{【写作提示】}#1\par}      % 本节该写什么
\newcommand{\material}[1]{\par\noindent\textbf{【可用素材】}#1\par}    % 已有的数据/结论清单
\newcommand{\pitfall}[1]{\par\noindent\textbf{【容易失分】}#1\par}     % 常见问题
\newcommand{\ph}[1]{\textbf{【#1】}}                                   % 行内占位符
\newcommand{\figph}[1]{%
  \begin{center}
  \fbox{\parbox{0.88\textwidth}{\centering\vspace{2ex}#1\vspace{2ex}}}
  \end{center}}

% 【新增】算法手册索引宏
\newcommand{\handbook}[2]{\par\noindent\textbf{【手册索引：#1】} #2\par}
\begin{document}

% ============================================================
% 【本文件使用说明】（注释，不进入 PDF）
% 一、文件由两部分组成：
%     A. 论文框架（正文区）：只有标题、表格骨架和“写作提示”，内容由你填写。
%     B. 原模板内容（文末）：原 1.tex 的全部正文，已逐行用 % 注释，未删除，
%        供你对照检查新框架有没有遗漏模板要求。
% 二、定位方法：
%     搜索“【写作提示】” -> 每个标题下该写什么、强调什么
%     搜索“【可用素材】” -> 你已算好、可直接引用的数字（怎么表述由你决定）
%     搜索“【容易失分】” -> 该节常见扣分点
%     搜索“【待填】”     -> 表格/占位需要填写的位置
% 三、写法建议：先按标题把空表、公式位置、每段主题句填出来，再回头润色语言。
% ============================================================

% ============================================================
% 【评分标准 → 本模板位置 对照表】（源自 5 篇国奖论文的行文规律）
% 评阅人实际按下面几件事给分，本模板已把它们对应到具体章节：
%   1) 摘要能否“一眼看到成果”      -> 摘要页：每问必须以“部分求解结果：某量 = 数值 单位”收尾
%   2) 是否真正读懂题目并转化        -> 二、问题分析：每问写“问题性质→关键中间量→方法→难点→处理手段”
%   3) 模型是否闭合、假设是否必要    -> 三、模型假设 + 五、模型搭建
%   4) 算法是否有针对性、步骤是否可复现 -> 六、模型求解：先讲清方法选择理由，再给离散格式与算法步骤
%   5) 结果是否可信                  -> 七、仿真验证：内部一致性 / 独立方法对照 / 物理与极限检验
%   6) 结论是否稳健                  -> 八、鲁棒性检验：参数区间+步长+曲线+变化百分比+敏感度排序
%   7) 是否有加分亮点                -> 各章“【可选亮点】”，以及给自创策略/算法命名
%   8) 格式与规范                    -> 附录（支撑材料、AI 声明、源程序）
% ============================================================

% ============================================================
% 【建模完善建议】（参考国奖论文的做法，逐条对应到本模板章节；写论文时按需取舍）
%  S1 无量纲化与特征数分析：把方程写成无量纲形式，用传质毕渥数、热扩散时间、
%     水分扩散时间等特征数解释“预热阶段先平衡、干燥由表及里”，
%     能同时支撑“一维简化”“两阶段划分”两个结论 -> 放入 5.2 节末尾。
%  S2 干燥速率曲线（DRC）与临界含水率：给出干燥速率随含水率的变化，识别
%     恒速段与降速段，与题目“预热平衡 + 恒温干燥”对应 -> 放入 6.8 节。
%  S3 收缩的两种处理对照：既按附件 2 直接插值半径，也从水分流失反推收缩率，
%     比较两种处理对干燥时长的影响 -> 放入 5.6 / 6.9 节。
%  S4 解析或渐近对照：扩散系数取常数时存在级数解，短时间存在渐近解，
%     可与数值解对照，作为最强的模型检验证据 -> 放入 7.2 节。
%  S5 守恒律检验：用导出数据验证“总水量变化 = 表面通量”的积分恒等式，
%     可精准定位柱坐标权重、边界半径等易错点 -> 放入 7.4 节（模板已留位置）。
%  S6 灵敏度协议化：参数区间 + 步长 + 曲线 + 变化百分比 + 敏感度排序，
%     并给出“控温控湿精度要求”的管理学结论 -> 放入第八章。
%  S7 优化推广：以“能耗/时间最小、品质达标”为目标、以烘房温湿度曲线为决策
%     变量，可把本文正问题升级为优化问题，体现模型的应用价值 -> 放入 9.3 节。
%  S8 命名关键方法：最大值判据、材料坐标变换、等效扩散放大因子、双方法互验、
%     积分守恒检验——每个名字至少出现两次（摘要 + 正文），强化记忆点。
% ============================================================

% ============================================================
% 【算法手册索引使用说明】
% 手册文件：A题_数学模型与算法手册.md（与 template.tex 同一目录）
% 正文中凡是算法相关的地方，都已插入“【手册索引：X.Y 小节名】”提示，
% 里面写明：读哪一节 + 读完后你至少要能写出哪几件事。
% 建议的阅读顺序（按写作进度读，不要一次全读）：
%   写“五、模型搭建”前   -> 手册 1.1 ~ 1.6；问题一~四分别再读 2.x / 3.x / 4.x / 5.3
%   写“六、模型求解”前   -> 手册 6.1 ~ 6.5（数值算法详解，每节含“为什么/怎么做”）
%   写“七、仿真验证”前   -> 手册 8.1 ~ 8.7（验证体系，含现成的残差表与结论句）
%   写“八、鲁棒性检验”前 -> 手册 8.3（收敛性做法）+ 8.6（误差量级）+ 9.3（脚本清单）
%   写论文语言卡壳时     -> 手册 1.8、1.9、8.7（现成段落模板）
% ============================================================

% ------------------------------------------------------------
% 【算法相关章节 ↔ 手册位置 速查表】
%   1.2/1.3 控制方程推导        -> 手册 1.1、1.2、1.3
%   1.4/1.5 边界与初始条件      -> 手册 1.4、1.5、1.6
%   1.7 无量纲化与特征数        -> 手册 1.7（可作 S1 亮点）
%   5.3~5.6 四问模型            -> 手册 2.1~2.3 / 3.1~3.2 / 4.1~4.2 / 5.1~5.3
%   6.2 空间离散（有限体积）    -> 手册 6.1
%   6.3 时间离散与稳定性        -> 手册 6.2
%   6.4 非线性与线性求解        -> 手册 6.4（非线性）+ 6.3（Thomas 追赶法）
%   6.5 后处理与精度控制        -> 手册 6.5
%   6.6~6.9 四问结果分析        -> 手册 2.4 / 3.3 / 4.3~4.4 / 5.4~5.5
%   7.2 双方法互验              -> 手册 8.2（+ 7.1、7.3 讲差异来源）
%   7.3 网格与时间步收敛        -> 手册 8.3
%   7.4 守恒恒等式检验          -> 手册 8.4
%   7.5 定性检验与误差分析      -> 手册 8.5、8.6
%   8.1/8.2 灵敏度分析          -> 手册 8.3（做法）+ 9.3（改脚本算扰动）
%   9.3 推广中的模型对比        -> 手册 5.2（三条路线）、7.3（差异来源）
% ------------------------------------------------------------

% ============================================================
% 摘要页（电子版第 1 页；页码从 1 开始；不写英文摘要；不得出现身份信息）
% ============================================================
\pagenumbering{arabic}
\setcounter{page}{1}

\begin{center}
{\zihao{-2}\bfseries \ph{论文标题：在原题基础上点出方法或亮点（例如“基于一维轴对称传热传质模型的……”）}}

\vspace{2ex}
{\zihao{3}\bfseries 摘\quad 要}
\end{center}

\prompt{\textbf{【本章评分点】} 国奖论文摘要的共同规律是：阅卷人只看这一页就能找到“模型名 + 方法名 + 每问的结果数字”。请严格按下面五段写，总长不超过一页。

第 1 段（总括，2$\sim$3 句）：现实背景 1 句 $\rightarrow$ “本文针对……问题，建立/运用……模型，基于……方法进行求解，解决了……问题” $\rightarrow$ 一句话点出全文主线索（例如“由守恒方程建模、用两种数值方法互相验证”）。\textbf{这一段必须出现模型名与方法名}。

第 2$\sim$4 段（每问一段，每段 4$\sim$6 句）：以“\textbf{针对问题$X$，}”开头，段内固定按“四步”推进：
（1）本问要求什么、与上一问的关系（继承或新增）；
（2）建模思路：先建什么量/什么方程，再用什么条件闭合（要出现关键量名，如特征长度、传质毕渥数、全场最大含水率、材料坐标）；
（3）关键处理或创新点，\textbf{并给它起一个名字}（例：“最大值判据”“材料坐标变换”“双方法互验”）；
（4）\textbf{段末必须写“部分求解结果：某量 = 数值 单位；某量 = 数值 单位，其余结果见附录/表 $X$”}——这是国奖论文最显著的特征，缺了就丢分。

第 5 段（收尾，2$\sim$3 句）：验证与鲁棒性各一句（必须带百分比或误差量级，例如“两种方法偏差小于 $0.4\%$”“参数扰动 $20\%$ 时干燥时长变化不超过 $X\%$”），最后一句写模型价值或可推广方向。

关键词 4$\sim$6 个，按“模型名 + 算法名 + 关键策略/物理量”排列，不要写“数学建模”这类空词。}

\handbook{1.8、1.9、8.7 现成段落模板}{写摘要时如果语言卡壳，可以先看手册 1.8（建模段落模板）、1.9（“模型建立”开头段）、8.7（“模型验证”段落范例），学它们的句式和信息密度，\textbf{但数字与结论必须换成你自己的}。}

\material{你已完成、可直接引用的数字（自己决定怎么表述，务必与正文表格一致）：
问题一：1800 s 中心温度 33.53 $^\circ$C、表面 36.76 $^\circ$C；表面含水率降至 1.51 kg/kg，中心仍为 2.55；
问题二：3 h 时径向温差仅约 0.07 $^\circ$C，含水率中心 1.77、表面 1.01；
问题三：干燥时长 57.15 h（约 2.38 天）；
问题四：考虑收缩时干燥时长 52.43 h；若忽略收缩，72 h 中心含水率仍有 0.22（无法达标）；
验证：问题一、三两种方法偏差 0.07\%$\sim$0.38\%；积分守恒律检验残差 $<0.5\%$；网格加密一倍误差减半。}

\pitfall{（1）全篇没有具体数字——国奖论文每一问都以“某量 = 数值 单位”收尾，缺数字等于放弃亮点；（2）只写“首先其次最后”的过程流水账，却不写算出了什么；（3）出现“结果很好、模型合理、效果显著”等无数据支撑的空话；（4）出现学校、姓名、赛区信息，或写了英文摘要、超过一页。}

% 【可选亮点】摘要里出现 1$\sim$2 个“自创名称”（如“最大值判据”“材料坐标变换”“双方法互验”“积分守恒检验”），能让评阅人立刻看出这篇论文有属于自己的工作。

\noindent{\heiti 关键词：}\ \ph{4$\sim$6 个，按“对象 + 方法 + 关键技术”排列，例如：传热传质耦合、一维轴对称、有限体积法、尺寸收缩、仿真验证、灵敏度分析}

\clearpage

% ============================================================
% 一、问题重述
% ============================================================
\section{问题重述}

\subsection{问题背景}

\prompt{用 5$\sim$8 行自己的语言概括：（1） 中药材干燥在产业中的地位与“能耗/品质”矛盾；（2） 热风干燥的两个阶段（预热平衡、恒温干燥）；（3） 影响干燥的关键因素（风温、湿度、风速、时长）；（4） 为什么需要机理建模与数值仿真（试验成本高、周期长）；（5） 本文针对本题要解决什么。}

\prompt{若引用文献，此处保留 1$\sim$3 篇高质量文献，并在正文标注引用编号；一句话说明文献与本文的关系（例如“已有研究给出了干燥动力学经验式，本文在此基础上建立可验证的传热传质模型”）。不要为凑数量引用无关文献。}

\subsection{问题重述}

\prompt{把题目翻译成建模语言，每个子问题压缩成 1 段或 1 行表格，统一采用“已知条件 $\rightarrow$ 决策/待求量 $\rightarrow$ 模型任务 $\rightarrow$ 输出要求”。\textbf{绝对不要大段照抄题面}。可以按下面的表格骨架填写（表中内容为提示，请自行替换为你的表述）：}

\begin{table}[H]
\centering
\caption{四个子问题的任务与输出要求}
\label{tab:restate}
\begin{tabular}{p{1.5cm}p{4.6cm}p{3.6cm}p{4.4cm}}
\toprule
问题 & 已知条件（题目给出什么） & 模型任务 & 输出要求\\
\midrule
1 & \ph{几何、初始条件、附件1环境数据、附录2物性} & \ph{建立什么模型} & \ph{表1、表2；result1.xlsx（1 s $\times$ 0.1 cm）}\\
\midrule
2 & \ph{附录3变物性、环境数据需外推} & \ph{建立什么模型} & \ph{表3、表4；result2.xlsx}\\
\midrule
3 & \ph{同问题2 + 达标要求} & \ph{求什么} & \ph{表5 + 干燥时长；result3.xlsx（60 s $\times$ 0.1 cm）}\\
\midrule
4 & \ph{附录4物性 + 附件2半径变化} & \ph{求什么} & \ph{表6（末列药材表面）；result4.xlsx}\\
\bottomrule
\end{tabular}
\end{table}

\pitfall{照抄题面超过半页；只写“问题 1 要求我们求温度”而不说清输入输出；忘记写清楚结果的保存形式（result*.xlsx、时间间隔、距离间隔、四位小数）。}

\clearpage

% ============================================================
% 二、问题分析
% ============================================================
\section{问题分析}

\prompt{本章只回答三件事：题目给了什么？真正的难点是什么？准备怎样把它数学化？\textbf{每个问题 4$\sim$10 行，不提前公布结果}，详细公式留给第五章。四问的写法建议如下。}

\subsection{总体分析思路}

\prompt{\textbf{【本章评分点】} 问题分析这一章决定评阅人对“你是否真读懂题”的第一印象。国奖论文的写法有两个共同点：
（1）每一问都用同一套句式推进：\textbf{“本问题是……（问题性质）。为求解它，首先需要确定……（关键中间量），可利用……（理论/方法）。由于……（难点或复杂度），因此在求解时采用……（处理手段）。”}
（2）分析里只出现“思路与依据”，不出现计算结果，也不展开具体公式。

本章先写一段总体思路：说明四问的共同本质与差异（例如：四问都由同一套守恒规律支配，差别只在物性公式、几何是否变化、输出目标不同），再给出技术路线：\textbf{物理机理 $\rightarrow$ 守恒方程 $\rightarrow$ 定解条件 $\rightarrow$ 离散求解 $\rightarrow$ 两种方法互验 $\rightarrow$ 灵敏度检验}。}

\handbook{阅读地图（手册开头）+ 9.1 各问写作要点}{手册最前面有一张“阅读地图”，用一条主线串起四问；9.1 给出四问的写作要点对照表（模型 / 关键公式 / 关键结论 / 对应结果文件）。写本章前先看这两处，可以保证你的“问题分析”与后面的“模型搭建、模型求解”口径一致。}

\figph{图 1\quad 技术路线图（占位）——建议画一条从左到右的主链：题目条件/附件数据 $\rightarrow$ 几何简化与假设 $\rightarrow$ 控制方程与边界条件 $\rightarrow$ 数值离散（有限体积/有限元）$\rightarrow$ 结果表格与图表 $\rightarrow$ 仿真验证与鲁棒性；下方用虚线引出“问题 1/2/3/4 的差异点在链条上的位置”。}

\subsection{问题一的分析}

\prompt{\textbf{针对问题一}，按“给了什么 $\rightarrow$ 难点在哪 $\rightarrow$ 怎样数学化”三步写：
（1）给了什么：规则圆柱几何、均匀初始场、附件 1 的离散环境数据、附录 2 的常数物性与扩散系数经验式；
（2）难点：\textbf{离散数据必须转成连续边界函数}，否则模型不闭合；扩散系数依赖未知的含水率（方程非线性）；圆柱坐标在轴心存在 $1/r$ 形式奇点；结果要在指定时刻与指定径向位置可核对；
（3）怎样数学化：先做量级判断（热扩散与水分扩散的特征长度与圆柱半长比较），据此把问题化归为\textbf{一维轴对称抛物型初边值问题}：内部导热 + 内部扩散，表面为对流（第三类）边界条件，轴心为对称条件；求解上用守恒型离散。
强调一句：“本问是后三问的基础模型，其方程与离散框架将被继承”。
\textbf{【句式模板】}“本问题要求在给定环境条件下求解温度场与水分浓度场。为求解它，首先需要把附件 1 的离散环境数据转化为连续边界函数，并把离散位置的含水率转化为连续场变量；可利用守恒定律建立控制方程。由于扩散系数随含水率变化、轴心存在坐标奇点，因此在离散时采用守恒型有限体积格式处理。”}

\pitfall{把“问题一就是求温度和水分”当分析；一上来就贴公式（分析章只需要说清思路与依据）；不说明为什么可以简化成一维。}

\paragraph{问题一分析（可直接写入正文）}
问题一的核心是确定时变烘房环境下药材内部温度和水分浓度的空间分布。
题目给出了规则圆柱几何、均匀初始状态、离散的烘房温度与水分浓度数据，
以及药材的热物性参数、对流传递参数和水分扩散系数经验式。因此，模型的
输入条件较为完整，但仍需解决离散边界数据连续化、圆柱坐标轴心奇异性以及
扩散系数非线性等问题。

首先，附件1只在若干离散时刻给出烘房温度和水分浓度，必须通过插值构造
连续边界函数，才能形成封闭的瞬态初边值问题。其次，药材可近似为细长圆柱，
且题目要求的结果均为径向位置上的数值。结合长径比和热、质扩散特征长度的
数量级比较，可将圆柱中截面的传热传质过程化归为一维轴对称径向问题。再次，
温度场由内部热传导和表面对流换热共同决定，水分场由内部扩散和表面对流传质
共同决定；其中扩散系数依赖未知水分浓度，使水分控制方程成为变系数非线性方程。

针对上述特点，本文由能量守恒、质量守恒以及Fourier和Fick定律建立温度场与
水分场控制方程，并分别施加轴心对称边界和表面对流边界。求解时采用守恒型
有限体积法处理圆柱坐标下的径向通量，在药材表面布置加密网格以分辨较大的
水分梯度，使用全隐式格式和Picard迭代处理时间推进与非线性系数更新。问题一
所得的控制方程和离散框架还将作为后续问题变物性、达标判据及收缩边界模型的基础。

\subsection{问题二的分析}

\prompt{\textbf{针对问题二}，重点写“继承 + 新增”：模型继承关系可写成一个箭头式表达（$M_1 \rightarrow M_1+\text{新增因素} \rightarrow M_2$），并回答三个问题：
（1）新增因素进入哪条方程/哪个系数？（物性 $\rho,c_p,k$ 与扩散系数 $D$ 变成含水率（和温度）的函数）
（2）这带来什么本质变化？（温度与水分由“互不影响”变成\textbf{双向耦合}：温度影响 $D$，水分影响物性与温度）
（3）算法上为什么要调整？（系数非线性，需要在每个时间步内迭代更新物性）
另外必须交代：干燥时长由 30 min 延长到数小时，附件 1 只覆盖 4 h，\textbf{需要明确外推规则}（恒温干燥阶段取末值）。
\textbf{【句式模板】}“本问题与问题一的差异只在物性：密度、比热容、导热系数与扩散系数不再是常数，而是自身状态的函数。为求解它，需要把物性写成耦合形式；由于温度与含水率互相影响，因此在每个时间步内都要重新计算物性并迭代，求解框架仍沿用问题一。”}

\pitfall{只写“问题二在问题一基础上增加了变物性”就结束——要写清“增加的量改变了哪条方程、为什么必须改算法”；忘记说明环境数据外推规则。}

\subsection{问题三的分析}

\prompt{\textbf{针对问题三}，核心是把“场”变成“判据”：题目要求“各处含水率低于 0.15”，这是一个\textbf{关于整个场的条件}。建议写三步：
（1）把要求数学化：定义全场最大值函数 $M(t)$，达标条件写成 $M(t)<0.15$，干燥时长定义为该条件首次成立的时刻；
（2）说明为什么可以只看中心：表面持续排湿、轴心为零通量面，含水率沿半径单调不减，故最大值在中心取得；再从“若内部出现局部极小则不会继续下降”的角度补一句极值原理的论证；
（3）说明怎么求：$C(0,t)$ 对时间单调，等价于求一个单调函数的零点，可用逐时步扫描或二分法。
强调：“本问的输出不是更多数据，而是一个时间数字”，因此判据的严谨性比场的精度更关键。
\textbf{【句式模板】}“本问题要求判断药材各部分是否达到质量要求，本质是一个关于整个场的条件，而不是新的一组数据。为求解它，需要把‘各处低于某值’转化为关于场最大值的条件；由于表面排湿、轴心绝湿，含水率沿半径单调下降，因此最大值只能出现在中心，判据可简化为中心值判据，问题化为求单调函数零点。”}

\pitfall{直接说“取中心含水率达到 0.15 就是干燥时间”而不给任何论证；不做单调性说明，导致判据看起来像假设。}

\subsection{问题四的分析}

\prompt{\textbf{针对问题四}，本问的“新东西”是\textbf{几何在变}，分析要围绕这一点展开：
（1）指出附件 2 给出的半径从 2.0 cm 单调减到 1.198 cm，意味着计算域本身在移动；
（2）说明不能照搬问题三：控制方程、边界条件、网格都必须随几何变化修正（边界通量中的半径要用当前半径）；
（3）给出解决思路：把物理坐标换成跟随材料收缩的\textbf{材料坐标}，使计算域重新固定；代价是方程中出现与 $R(t)$ 有关的缩放因子，物理上表现为“扩散路程变短、单位干物质表面积增大”，两者都使干燥加快；
（4）预告对照设计：再做一次“忽略收缩”的计算，用定量对比说明尺寸变化的必要性（这是本问最容易被评阅人认可的亮点）。
\textbf{【句式模板】}“本问题与前三问的本质差别在于几何本身在变化。为求解它，需要把物理坐标换成随材料一起收缩的坐标，使计算域重新固定；由于坐标变换会在方程中引入与当前半径有关的缩放因子，因此在离散时每个时间步都要更新几何系数与表面通量中的半径，并额外计算一个不收缩的对照模型以量化尺寸变化的影响。”}

\pitfall{只说“半径变小所以干得更快”而不说明如何处理动边界；不写“距离列到底指初始位置还是当前位置”这个坐标约定，导致表 6 被误解。}

\clearpage

% ============================================================
% 三、模型假设
% ============================================================
\section{模型假设}

\prompt{\textbf{【本章评分点】} 假设不是越多越好，而是“必要、合理、可解释、可检验”。国奖论文的共同写法是：\textbf{每条假设 = 一句假设 + 一句理由}，理由优先用\textbf{量级比较}或\textbf{题目条件}，绝不用“为了简化计算”。假设尽量控制在 5$\sim$8 条，并且后面每一章都要能对上它（第七章误差分析、第八章假设放宽都要回头引用这里的条目）。

下面列出本题应覆盖的 8 个方面，请用你自己的语言写成条目（示例见文末归档区原模板）。}

\handbook{0.3 基本假设}{手册 0.3 给了 8 条假设及其理由（含端面占比 $R/L=8\%$、热扩散时间约 $39$ min 等具体依据）。读完后你要能自己写出：（1） 每条假设的一句话表述；（2） 支撑它的量级或题目条件；（3） 它后面会在哪一章被检验。不要照抄手册的理由，改成你自己的量级估算。}

\begin{enumerate}[label=(\arabic*)]
  \item \ph{连续介质与局部热平衡}：\ph{理由（提示：用热扩散时间的量级说明）}
  \item \ph{初始温度与含水率均匀分布}
  \item \ph{一维轴对称（忽略轴向梯度与端面）}：\ph{理由（提示：长径比、特征长度与半长的比较）}
  \item \ph{水分迁移服从 Fick 扩散，扩散系数用题目经验式}
  \item \ph{表面按牛顿冷却与对流传质交换热量和水分}
  \item \ph{内部无体热源、无水分源，忽略辐射与交叉效应}
  \item \ph{物性只依赖当前状态（不依赖历史）}
  \item \ph{问题四的收缩方式假设（提示：均匀/同比例收缩）}
\end{enumerate}

\prompt{建议再补一张“关键假设影响表”：假设 $\rightarrow$ 若放宽会怎样 $\rightarrow$ 对结论的影响量级。这张表能把“假设”与第七章的误差分析连起来，很受评阅人欢迎。}

\begin{table}[H]
\centering
\caption{关键假设的影响分析（骨架，请填写）}
\label{tab:assumption-check}
\begin{tabular}{p{3.2cm}p{5.2cm}p{6.2cm}}
\toprule
假设 & 若放宽会怎样 & 对结论的影响/可检验性\\
\midrule
\ph{一维轴对称} & \ph{需二维/三维计算} & \ph{端面面积占比 $R/L=8\%$，可在误差分析中说明}\\
\ph{忽略蒸发潜热} & \ph{升温会滞后} & \ph{题目未给潜热等参数，属参数不可辨识的简化}\\
\ph{均匀收缩} & \ph{非均匀收缩会使结果略变} & \ph{可与仿真位移场对比说明}\\
\ph{其他} & & \\
\bottomrule
\end{tabular}
\end{table}

\pitfall{写“假设模型结果正确”“假设数据准确”（循环论证）；假设数量很多但每条都没有理由；假设与后面的模型、误差分析互不呼应。}

\clearpage

% ============================================================
% 四、符号说明
% ============================================================
\section{符号说明}

\prompt{只列正文反复使用的核心符号；单位必须写清楚；同一符号全文不得一符多义。建议分成四组列：空间与时间、场变量、物性与系数、数值与判据。临时求和指标不必列入。}

\handbook{0.2 符号表}{手册 0.2 是本题的完整符号表（含 $X$、$R_t$、$M(t)$、$t_{dry}$、$V_i$ 等本题特有符号）。可比对取舍，但要保证\textbf{符号表与正文实际出现的符号一致}：正文没用到的不列，用到的一定要在表里。}

\begin{table}[H]
\centering
\caption{主要符号说明（骨架，请填写）}
\label{tab:symbol}
\begin{tabular}{p{2.4cm}p{9.0cm}p{2.6cm}}
\toprule
符号 & 含义 & 单位\\
\midrule
\ph{$r$} & \ph{到药材中心轴线的径向距离} & \ph{m}\\
\ph{$t$} & \ph{时间} & \ph{s}\\
\ph{$T$、$C$} & \ph{药材温度、干基含水率} & \ph{$^\circ$C、kg/kg}\\
\ph{$T_a$、$C_a$} & \ph{烘房温度、烘房含水率（附件 1）} & \ph{$^\circ$C、kg/kg}\\
\ph{$R_0$、$R(t)$} & \ph{初始半径、问题四中随时间变化的半径} & \ph{m}\\
\ph{$\rho$、$c_p$、$k$} & \ph{密度、比热容、导热系数（附录公式）} & \ph{kg/m$^3$、J/(kg$\cdot$K)、W/(m$\cdot$K)}\\
\ph{$D$} & \ph{水分扩散系数（附录公式）} & \ph{m$^2$/s}\\
\ph{$h$、$h_m$} & \ph{对流换热系数、对流传质系数} & \ph{W/(m$^2\cdot$K)、m/s}\\
\ph{$M(t)$、$t_{dry}$} & \ph{全场最大含水率、干燥时长} & \ph{kg/kg、h}\\
\ph{$N$、$\Delta r$、$\Delta t$} & \ph{控制体数目、网格步长、时间步长} & \ph{---、m、s}\\
\ph{其他} & & \\
\bottomrule
\end{tabular}
\end{table}

\pitfall{符号表里塞很多只有一个公式用到的中间量；符号与正文不一致（例如表里叫 $C_m$、正文写 $C_{air}$）；忘记写单位。}

\clearpage

% ============================================================
% 五、模型搭建
% ============================================================
\section{模型搭建}

\prompt{\textbf{【本章评分点】} 这是全文最重要的一章，评阅人看三件事：\textbf{（1）模型是否闭合}（变量、方程、初边值条件是否齐全，能不能解）；\textbf{（2）每一步是否有出处}（来自物理规律还是题目条件）；\textbf{（3）四问之间是否有清晰的继承关系}（而不是四套互不相关的模型）。

只做一件事：把题目条件一步步变成\textbf{封闭}的数学定解问题。建议按“几何简化 $\rightarrow$ 守恒方程 $\rightarrow$ 边界与初始条件 $\rightarrow$ 物性闭合 $\rightarrow$ 完整定解问题”的顺序组织；公共部分只写一次（5.2 节），四问的差别分别写在 5.3$\sim$5.6 节。

\textbf{【可选亮点】} 给关键处理手段起名字，并在正文中始终用这个名字（例如“最大值判据”“材料坐标变换”“等效扩散放大因子”）。国奖论文普遍这么写（如“无效点判断策略”“分区域同心圆规划策略”），评阅人能据此迅速记住你的工作。}

\subsection{建模总路线}

\prompt{用一张“文字到公式”的路线图或一段话说明总思路（题目条件 $\rightarrow$ 变量定义 $\rightarrow$ 基本规律 $\rightarrow$ 数学关系 $\rightarrow$ 定解条件 $\rightarrow$ 完整模型）。不需要新公式。}

\subsection{公共物理模型}

\prompt{本节把四问共用的部分写全，建议分四小块：
（1）\textbf{几何简化}：给出长径比等几何量，并用特征长度（热扩散特征长度、水分扩散特征长度）与圆柱半长做量级比较，据此说明“可忽略轴向与端面、化为一维轴对称问题”。\textbf{这一步的“为什么”比公式本身更重要}；
（2）\textbf{控制方程}：由“圆柱薄壳微元的守恒分析”分别导出水分浓度方程（Fick 扩散 + 质量守恒）与温度方程（Fourier 导热 + 能量守恒）。提示：两条方程可以用同一套微元法写一次、分别用不同守恒量即可，篇幅省一半；
（3）\textbf{边界与初始条件}：表面为对流（第三类）边界条件（温度用换热系数、水分用传质系数），轴心为对称（零通量）条件，初始条件为均匀温度与均匀含水率；
（4）\textbf{环境边界函数}：附件 1 为离散数据，需构造连续函数；说明为何选\textbf{分段线性插值}（不产生非物理过冲），并明确超过数据范围后的\textbf{外推规则}（恒温干燥阶段取末值）。}
\handbook{1.1$\sim$1.6 控制方程推导}{手册第 1 章用“圆柱薄壳微元”一次性把两条方程推完（1.1 取微元、1.2 水分守恒、1.3 能量守恒、1.4 边界条件、1.5 初始条件、1.6 完整定解问题）。读完后你要能自己写出：（1）微元的内外表面积与体积；（2）“流入 $-$ 流出 $=$ 累积”这一步如何变成微分方程；（3）两条方程里物性系数分别出现在什么位置；（4）表面为什么用第三类条件、轴心为什么是零通量。\textbf{建议先在纸上按手册重推一遍，再动笔写正文}，公式请自行排版、统一编号并检查量纲。}

\handbook{1.7 无量纲化与特征数（可选亮点）}{手册 1.7 给出特征数与四问取值表（传质毕渥数、热扩散时间、水分扩散时间）。加上这一小节，可以用“热扩散时间远小于水分扩散时间”解释“预热阶段先平衡”，用“传质毕渥数与 1 同量级”解释“表面与内部阻力都要建模”——这是把“会算”升级为“能解释”的最省力做法。}

% 【可复制素材】控制方程、边界条件、特征长度定义见 A题_数学模型与算法手册.md 第 1 章（式 1~式 7）。

\subsection{问题一的模型}

\prompt{\textbf{针对问题一}：写出附录 2 的物性（密度、比热容、导热系数为常数，扩散系数为含水率的经验函数），并说明“本问温度方程不含含水率、水分方程不含温度，故两个子模型不显式耦合，可分别求解”。再补一句物性趋势分析（扩散系数随含水率增大而增大，故最大值出现在初始时刻），为第六章解释“只有表层失水”做铺垫。}

\handbook{2.1 / 2.2 / 2.3 问题一的物性、环境插值、定解问题}{手册 2.2 讲了“为什么用分段线性插值、超范围为什么取常数外推”，2.3 把问题一的定解问题写完整。读完后你要能写出：（1） 附录 2 的三个常数物性与扩散系数经验式；（2） 环境函数的两条插值/外推规则；（3） 为什么本问两条方程不耦合、可以分别求解。}

\paragraph{问题一模型建立（可直接写入正文）}
将药材近似为半径 $R=0.02\ \mathrm{m}$、长度 $L=0.25\ \mathrm{m}$ 的均匀圆柱体。
题目中的“到药材中心距离”解释为圆柱中截面上到中心轴线的径向距离。药材长径比为
\begin{equation}
\frac{L}{2R}=6.25.
\end{equation}
问题一中的热扩散率为
\begin{equation}
\alpha=\frac{k}{\rho c_p}
 =\frac{0.36}{820\times2600}
 =1.6886\times10^{-7}\ \mathrm{m^2/s}.
\end{equation}
在 $t=1800\ \mathrm{s}$ 时，热扩散特征长度为
\begin{equation}
\ell_T=\sqrt{\alpha t}\approx1.743\ \mathrm{cm}.
\end{equation}
水分扩散系数为
\begin{equation}
D(C)=7\times10^{-9}\exp\left(-\frac{0.89}{C}\right).
\end{equation}
由
\begin{equation}
\frac{\mathrm dD}{\mathrm dC}=D(C)\frac{0.89}{C^2}>0
\end{equation}
可知，$D(C)$ 随含水率增大而增大。以初始含水率
$C_0=2.55\ \mathrm{kg/kg}$ 计算，扩散系数上界为
\begin{equation}
D_{\max}=4.9377\times10^{-9}\ \mathrm{m^2/s},
\end{equation}
相应的水分扩散特征长度为
\begin{equation}
\ell_C=\sqrt{D_{\max}t}\approx0.298\ \mathrm{cm}.
\end{equation}
由于 $\ell_T$ 和 $\ell_C$ 均远小于圆柱半长 $L/2=12.5\ \mathrm{cm}$，
端面扰动在问题一的计算时间内难以传播至中截面。因此忽略轴向和周向梯度，
将模型简化为 $0\le r\le R$ 上的一维轴对称径向问题。

附件1给出了离散时刻的烘房温度和水分浓度。对于
$t_i\le t\le t_{i+1}$，采用分段线性插值构造连续边界函数：
\begin{equation}
T_a(t)=T_{a,i}+\frac{T_{a,i+1}-T_{a,i}}{t_{i+1}-t_i}(t-t_i),
\end{equation}
\begin{equation}
C_a(t)=C_{a,i}+\frac{C_{a,i+1}-C_{a,i}}{t_{i+1}-t_i}(t-t_i).
\end{equation}

根据能量守恒和Fourier热传导定律，药材内部温度满足
\begin{equation}
\boxed{
\rho c_p\frac{\partial T}{\partial t}
=\frac{1}{r}\frac{\partial}{\partial r}
\left(kr\frac{\partial T}{\partial r}\right)},
\qquad 0<r<R,\quad 0<t\le1800.
\label{eq:merged-q1-heat}
\end{equation}
根据质量守恒和Fick扩散定律，药材内部水分浓度满足
\begin{equation}
\boxed{
\frac{\partial C}{\partial t}
=\frac{1}{r}\frac{\partial}{\partial r}
\left[rD(C)\frac{\partial C}{\partial r}\right]},
\qquad 0<r<R,\quad 0<t\le1800.
\label{eq:merged-q1-moisture}
\end{equation}

烘干开始时药材内部状态均匀分布，初始条件为
\begin{equation}
T(r,0)=28,\qquad C(r,0)=2.55,
\qquad 0\le r\le R.
\label{eq:merged-q1-initial}
\end{equation}
在中心轴线 $r=0$ 处，由轴对称性有
\begin{equation}
\left.\frac{\partial T}{\partial r}\right|_{r=0}=0,
\qquad
\left.\frac{\partial C}{\partial r}\right|_{r=0}=0.
\label{eq:merged-q1-center}
\end{equation}
在药材表面 $r=R$ 处采用第三类边界条件：
\begin{equation}
-k\left.\frac{\partial T}{\partial r}\right|_{r=R}
=h_T\left[T(R,t)-T_a(t)\right],
\label{eq:merged-q1-heat-bc}
\end{equation}
\begin{equation}
-D(C)\left.\frac{\partial C}{\partial r}\right|_{r=R}
=h_C\left[C(R,t)-C_a(t)\right].
\label{eq:merged-q1-moisture-bc}
\end{equation}
其中 $h_T=25\ \mathrm{W/(m^2\cdot K)}$，$h_C=8\times10^{-7}\ \mathrm{m/s}$。
由于问题一中温度方程不含水分浓度，水分方程也不含温度，两个子模型不发生
显式耦合，可以分别求解。问题一不显式考虑蒸发潜热反馈，原因是题目未给出
汽化潜热、干物质密度及蒸发区域等必要参数。

\subsection{问题二的模型}

\prompt{\textbf{针对问题二}：写出附录 3 的四条经验公式（密度、比热容、导热系数、扩散系数），特别说明扩散系数指数项中的温度必须是\textbf{开尔文}。然后说明耦合结构：含水率下降 $\rightarrow$ 物性改变 $\rightarrow$ 温度分布改变 $\rightarrow$ 扩散系数改变 $\rightarrow$ 反过来影响含水率。除物性外，方程与边界条件与问题一完全相同，这就是“模型继承”的具体含义。}

\handbook{3.1 / 3.2 问题二的物性与两阶段刻画}{手册 3.1 列出附录 3 的四条经验式并强调温度取开尔文；3.2 说明“预热平衡 + 恒温干燥”两个阶段如何仅由环境函数体现。读完后你要能写出：（1） 四条物性公式；（2） 耦合链条（含水率 $\rightarrow$ 物性 $\rightarrow$ 温度 $\rightarrow$ 扩散系数 $\rightarrow$ 含水率）；（3） 为什么不必人为给两个阶段划分界线。}

\subsection{问题三的模型}

\prompt{\textbf{针对问题三}：不必重复写方程，只写“新增的判据”：
（1）定义全场最大含水率 $M(t)$ 与干燥时长 $t_{dry}$（“首次满足 $M(t)<0.15$ 的时刻”）；
（2）说明判据可简化为中心判据的理由（表面排湿、轴心零通量、沿半径单调不减；并说明内部局部极小不会先达标）；
（3）说明求解方式：单调函数的零点问题（逐时步扫描或二分法），并写出所选方法的停机准则。}

\handbook{4.1 / 4.2 判据的数学化与“只看中心”的论证}{手册 4.1 给出最大值判据的定义式，4.2 给出“为什么可以只看中心”的两步论证（沿半径单调不减 + 局部极小不会先达标）。这两段是问题三最容易被追问的地方，\textbf{建议按手册的论证层次自己重写一遍}，不要只写结论。}

\subsection{问题四的模型}

\prompt{\textbf{针对问题四}：本问差异全部集中在“几何在动”，建议按五步写：
（1）\textbf{收缩数据处理}：附件 2 给出半径随时间变化，用同样的分段线性插值构造半径函数，并给出首末值；
（2）\textbf{收缩运动学假设}：写出均匀收缩的坐标映射（材料点初始位置与当前半径的关系），并说明它与“变形域 + 边界位移”的位移场一致；
（3）\textbf{方程变换}：把 $r$ 换成跟随材料的坐标，重新做一次守恒分析，得到带缩放因子的水分与温度方程（提示：面积因子与梯度变换中的比例因子正好抵消，这是选材料坐标的关键好处，值得写一句）；
（4）\textbf{边界条件变换}：说明表面通量中的半径要换成\textbf{当前半径}，并给出变换后的表达式；
（5）\textbf{对照模型}：说明还要计算“忽略收缩”的模型，用于第六章定量比较。
提示：材料坐标下的表达式与等价形式在《A题\_数学模型与算法手册》第 5 章已推导完整。另外务必声明\textbf{表 6 中“到药材中心的距离”指材料坐标（初始位置）还是当前几何位置}。}

\handbook{5.1$\sim$5.3、5.6 问题四的收缩建模（本文最难的算法段）}{手册 5.1 讲半径数据的处理；5.2 列出三条可选路线（忽略收缩 / 材料坐标 / 归一化坐标）并说明为什么选材料坐标；5.3 用五步给出完整推导（运动学假设 $\rightarrow$ 守恒分析 $\rightarrow$ 方程变换 $\rightarrow$ 边界条件变换 $\rightarrow$ 等价形式）；5.6 给出它与 COMSOL“变形几何”操作的对应关系。\textbf{读完后你必须能自己写出这五步}，否则这一段会成为全文最薄弱的环节。建议按“手册 5.3.1 $\to$ 5.3.2 $\to$ 5.3.3”逐段读、每读完一段就合上手册复述一遍。}

\subsection{模型汇总}

\prompt{用一张表把四问并列起来（几何、物性来源、是否耦合、关键修正、核心公式编号），让评阅人一眼看出四问的继承关系与差异。}

\begin{table}[H]
\centering
\caption{四个子问题的数学模型汇总（骨架，请填写）}
\label{tab:model-summary}
\begin{tabular}{p{1.5cm}p{3.3cm}p{2.8cm}p{3.3cm}p{3.4cm}}
\toprule
问题 & 控制方程 & 物性来源 & 关键修正 & 核心公式编号\\
\midrule
1 & \ph{式(\ )} & \ph{附录 2} & \ph{无耦合} & \ph{式(\ )}\\
2 & \ph{式(\ )} & \ph{附录 3} & \ph{变物性双向耦合、环境外推} & \ph{式(\ )}\\
3 & \ph{同问题 2} & \ph{附录 3} & \ph{达标判据} & \ph{式(\ )}\\
4 & \ph{式(\ )} & \ph{附录 4} & \ph{材料坐标 + 动边界} & \ph{式(\ )}\\
\bottomrule
\end{tabular}
\end{table}

\pitfall{方程推导没有任何出处（既没物理规律也没微元图）；只列公式不写边界条件导致模型不封闭；四问分别重写一遍同样的方程，篇幅浪费且看不出继承关系；漏写问题四的坐标约定。}

\clearpage

% ============================================================
% 六、模型求解
% ============================================================
\section{模型求解}

\prompt{\textbf{【本章评分点】} 评阅人关心两件事：\textbf{算法是否与模型匹配}（为什么用这个方法、关键步骤能不能照着复现）、\textbf{结果是否被读懂}（有没有解释趋势与物理原因）。国奖论文的固定套路是“\textbf{先讲方法选择理由，再给离散格式与算法步骤，最后给结果表 + 一段结果分析}”，结果分析绝不重复念数字。}

\subsection{求解总体方案}

\prompt{用 3$\sim$5 行 + 一张流程图说明总体流程：读入附件数据并构造环境/半径函数 $\rightarrow$ 生成径向控制体 $\rightarrow$ 时间推进（更新物性 $\rightarrow$ 组装三对角方程组 $\rightarrow$ 求解）$\rightarrow$ 抽取到题目要求的位置与时刻 $\rightarrow$ 输出 Excel。强调“内部计算网格/时间步”与“题目要求的输出位置/时刻”是两回事。}

\handbook{第 6 章总览（6.1$\sim$6.5）}{手册第 6 章把本文用到的数值算法逐个讲清（每个算法都按“是什么 / 为什么用它 / 怎么做”三段写）：6.1 有限体积法、6.2 全隐式时间离散与稳定性、6.3 Thomas 追赶法、6.4 非线性物性处理、6.5 后处理。\textbf{建议先把 6.1$\sim$6.3 读完再动笔写本章}，这样你能自己判断“哪一步是必须写进论文的、哪一步只是实现细节”。}

\figph{图 2\quad 求解流程图（占位）——建议体现“两条并行线：自主有限体积程序”与“COMSOL 有限元仿真”，最后汇入“结果对比与验证”。}

\subsection{问题一的自主数值求解}

为求解式\eqref{eq:merged-q1-heat}--式\eqref{eq:merged-q1-moisture-bc}，
本文采用一维轴对称有限体积法进行空间离散，并采用全隐式格式推进时间。
题目规定的 $0.1\ \mathrm{cm}$ 为空间输出间隔，不直接作为内部计算网格。
考虑到药材表面附近可能形成较大的水分梯度，内部径向网格向 $r=R$ 逐渐加密，
计算完成后再将结果抽取到题目要求的位置。

设径向节点为
\begin{equation}
0=r_0<r_1<\cdots<r_N=R,
\end{equation}
节点 $r_i$ 对应控制体的内、外表面分别为 $r_{i-\frac12}$ 和
$r_{i+\frac12}$。忽略公共圆周因子后，第 $i$ 个控制体的几何权重为
\begin{equation}
V_i=\frac12\left(r_{i+\frac12}^2-r_{i-\frac12}^2\right).
\label{eq:merged-q1-volume}
\end{equation}

温度方程的全隐式离散形式为
\begin{equation}
\begin{aligned}
\rho c_pV_i\frac{T_i^{n+1}-T_i^n}{\Delta t}
={}&kr_{i+\frac12}
\frac{T_{i+1}^{n+1}-T_i^{n+1}}{r_{i+1}-r_i}\\
&-kr_{i-\frac12}
\frac{T_i^{n+1}-T_{i-1}^{n+1}}{r_i-r_{i-1}}.
\end{aligned}
\label{eq:merged-q1-heat-fvm}
\end{equation}

水分浓度方程的全隐式离散形式为
\begin{equation}
\begin{aligned}
V_i\frac{C_i^{n+1}-C_i^n}{\Delta t}
={}&r_{i+\frac12}D_{i+\frac12}^{n+1}
\frac{C_{i+1}^{n+1}-C_i^{n+1}}{r_{i+1}-r_i}\\
&-r_{i-\frac12}D_{i-\frac12}^{n+1}
\frac{C_i^{n+1}-C_{i-1}^{n+1}}{r_i-r_{i-1}}.
\end{aligned}
\label{eq:merged-q1-moisture-fvm}
\end{equation}

相邻控制体界面的扩散系数取调和平均：
\begin{equation}
D_{i+\frac12}=\frac{2D_iD_{i+1}}{D_i+D_{i+1}}.
\label{eq:merged-q1-harmonic}
\end{equation}
该处理将相邻控制体的扩散阻力串联起来，能够保持界面水分通量连续。

在轴线 $r=0$ 处，控制体内侧表面积为零，内侧通量自然消失，无需直接计算
连续方程中的 $1/r$ 项。在最外层控制体中，将外侧通量分别替换为
\begin{equation}
Rh_T\left[T_a(t^{n+1})-T_N^{n+1}\right]
\end{equation}
和
\begin{equation}
Rh_C\left[C_a(t^{n+1})-C_N^{n+1}\right].
\end{equation}

温度方程为线性三对角方程组，每个时间步求解一次即可。水分方程中的扩散系数
$D(C)$ 依赖新时刻的未知水分浓度，因此采用Picard迭代：以 $C^n$ 作为初始迭代值，
根据当前迭代结果更新扩散系数，再求解新的水分场。当
\begin{equation}
\max_i\left|C_i^{n+1,(m+1)}-C_i^{n+1,(m)}\right|<\varepsilon_C
\end{equation}
时，认为当前时间步收敛。两个方程组均采用Thomas追赶法求解，单次求解复杂度为
\begin{equation}
O(N).
\end{equation}

内部时间步 $\Delta t_{\mathrm{in}}$ 与题目要求的输出步长
$\Delta t_{\mathrm{out}}=1\ \mathrm{s}$ 相互独立。程序采用较小内部时间步推进，
仅在整数秒保存结果；空间上则将内部解抽取至
\begin{equation}
r=0,0.001,\ldots,0.020\ \mathrm{m}
\end{equation}
的位置。进一步增加内部径向节点数并减小内部时间步，若题目指定位置的结果变化
低于允许误差，则认为数值解达到空间和时间收敛。

\subsection{空间离散}

…% %   \item 第1页为承诺书，第2页为编号专用页。
% %   \item 第3页为摘要专用页；摘要含标题和关键词，原则上不超过1页，不需要英文摘要。
% %   \item 从摘要页开始编写页码，页码位于页脚中部，以阿拉伯数字1开始连续编号。
% %   \item 第4页开始为正文；\textbf{不要目录}；正文不超过30页。
% %   \item 正文之后为附录，附录页数不限；附录必须打印并与正文装订提交。
% %   \item 摘要、正文和附录任何地方都不能显示参赛者身份、所在学校及赛区信息。
% %   \item 所有引用他人或公开资料（包括网上资料）的成果必须按科技论文规范列出参考文献，并在正文引用处标注。
% %   \item 未规定的字号、字体、行距、颜色等不统一要求；如赛区另有规定，应同时遵守赛区要求。
% % \end{enumerate}
% 
% % \subsection*{电子版}
% % \begin{enumerate}
% %   \item 参赛队提交“参赛论文”和“支撑材料”两个电子文件。
% %   \item 参赛论文必须是一个单独文件，PDF或Word二选一，建议PDF；文件大小不超过20MB。
% %   \item 电子版论文必须与纸质版内容及格式（包括附录）完全一致。
% %   \item \textbf{承诺书和编号专用页不能放入电子版论文}；因此电子版第一页必须是摘要专用页。
% %   \item 支撑材料统一压缩为一个RAR或ZIP文件，大小不超过20MB。
% %   \item 支撑材料文件列表必须写入论文附录。
% %   \item 如果没有支撑材料，可以不提供支撑材料文件，但附录必须注明“本论文没有支撑材料”。
% %   \item 承诺书和编号专用页不要放入支撑材料；所有文件中均不得显示参赛者身份、学校及赛区信息。
% % \end{enumerate}
% 
% % \subsection*{附录与资格风险}
% % \begin{itemize}
% %   \item 缺少必要源程序；
% %   \item 程序无法运行；
% %   \item 程序运行结果与论文不一致；
% %   \item 支撑材料与论文内容不相符；
% %   \item 违反身份信息隐匿要求；
% %   \item 不按要求提交AI声明或AI详情材料；
% % \end{itemize}
% % \textbf{均可能带来严重后果，甚至影响评奖资格。}
% 
% % ============================================================
% % 十七、最终交稿检查清单
% % ============================================================
% \section{最终交稿：30分钟“红线检查”}
% 
% \begin{longtable}{p{1cm}p{2cm}p{12cm}}
% \toprule
% 序号 & 检查项 & 检查标准 \\
% \midrule
% 1 & 身份信息 & 摘要、正文、附录、支撑材料均无姓名、学校、赛区等身份信息 \\
% 2 & 页数 & 正文不超过30页；附录不限 \\
% 3 & 目录 & 正文前没有目录 \\
% 4 & 页码 & 从摘要页开始为1，页脚中部连续编号 \\
% 5 & 摘要 & 原则上不超过1页；包含模型、方法、主要结果、结论 \\
% 6 & 题目重述 & 没有大段复制原题 \\
% 7 & 模型假设 & 每个关键假设都有实际理由 \\
% 8 & 符号 & 符号统一、首次出现有解释、单位明确 \\
% 9 & 模型 & 每个核心问题都有明确数学模型 \\
% 10 & 约束 & 优化模型的每条关键约束都有现实来源 \\
% 11 & 算法 & 选择理由、关键步骤、改进点写清楚 \\
% 12 & 图表 & 编号、标题、单位、图例、数据均正确 \\
% 13 & 结果 & 正文真正解释了最关键的结果 \\
% 14 & 检验 & 至少有与问题性质匹配的验证 \\
% 15 & 灵敏度 & 关键参数变化不会导致结论无故失效，或已说明影响 \\
% 16 & 评价 & 优缺点与前文假设/检验呼应 \\
% 17 & 文献 & 每篇引用都在正文出现；每处重要外部资料都有来源 \\
% 18 & AI声明 & 按实际情况二选一，放在参考文献之前 \\
% 19 & AI详情 & 使用AI时有“AI工具使用详情.pdf” \\
% 20 & 源程序 & 全部完整、可运行、结果与论文一致 \\
% 21 & 支撑材料 & 文件列表写入附录，RAR/ZIP不超过20MB \\
% 22 & 电子版 & 单独PDF/Word，不含承诺书和编号专用页 \\
% 23 & 纸质版 & A4、页边距至少2.5cm、左侧装订 \\
% 24 & 最终复现 & 从零开始运行程序能复现论文关键表格和图形 \\
% \bottomrule
% \end{longtable}
% 
% % ============================================================
% % 十八、培训最后强调的“论文高分逻辑”
% % ============================================================
% \section{培训总结：把论文写成一条完整的证据链}
% 
%  
% 
% \tip{真正高质量的数模论文不是“公式最多”“算法最复杂”“图最多”，而是每一页都在推进证据链。评阅人能够从题目条件一路追踪到最终结论，并且能够相信：假设合理、模型闭合、算法有效、结果可信、结论可复现。}
% 
% \begin{table}[H]
% \centering
% \caption{一篇优秀数模论文的五个关键词}
% \begin{tabular}{p{3cm}p{10cm}}
% \toprule
% 关键词 & 具体含义 \\
% \midrule
% 清楚 & 变量、模型、算法、结果一眼可读 \\
% 严谨 & 条件、假设、公式、数据、引用均有依据 \\
% 有效 & 模型真正回答题目，而不是为了套算法 \\
% 可信 & 有检验、误差、灵敏度或稳健性分析 \\
% 可复现 & 程序、数据和关键步骤足以复现核心结果 \\
% \bottomrule
% \end{tabular}
% \end{table}
% 
% \begin{center}
% \vspace{2ex}
% \Large\bfseries
% 不是把“做过的东西”全部写进论文，\par
% 而是把“最有说服力的证据”组织成一条完整的逻辑链。
% \end{center}
% 
% \end{document}
% 

\end{document}
